\documentclass[a11paper]{article}

\usepackage{tabularx}
\usepackage{libs/titlepage}
\usepackage{libs/document}
\usepackage{booktabs}
\usepackage{multicol}
\usepackage{float}
\usepackage{varwidth}
\usepackage{tcolorbox}
\usepackage{tikz}
\usepackage{graphicx}
\usepackage{pifont}
\usepackage{framed}
% \usepackage[toc,page]{appendix}
\usepackage[dvipsnames]{xcolor}

\title{Observations clés issues de la visite chez Polytech}

\class{Innovation et création de produits}
\classnb{GEN490}

\teacher{Joao Gabriel Alves Ribeiro Rosa}

\author{
  \addtolength{\tabcolsep}{-0.4em}
  \begin{tabular}{rcl} % Ajouter des auteurs au besoin
      Benjamin Chausse & -- & CHAB1704 \\
  \end{tabular}
}

\newcommand{\todo}[1]{\begin{color}{Red}\textbf{TODO:} #1\end{color}}
\newcommand{\note}[1]{\begin{color}{Orange}\textbf{NOTE:} #1\end{color}}
\newcommand{\fixme}[1]{\begin{color}{Fuchsia}\textbf{FIXME:} #1\end{color}}
\newcommand{\question}[1]{\begin{color}{ForestGreen}\textbf{QUESTION:} #1\end{color}}

% Define the command
\newcommand{\boxfig}[3]{%
  \begin{figure}[H]
    \centering
    \begin{tcolorbox}[colback=gray!5!white, colframe=black!50, boxrule=0.5pt,
                      width=0.9\linewidth, arc=3pt, auto outer arc,
                      boxsep=5pt, left=5pt, right=5pt, top=5pt, bottom=5pt]
      #3
    \end{tcolorbox}
    \caption{#1}
    \label{box:#2}
  \end{figure}
}

\newcommand{\quickfig}[4]{%
  \begin{figure}[H]
    \centering
    \includegraphics[width=#3\textwidth]{#4}
    \caption{#1}
    \label{fig:#2}
  \end{figure}
}

\begin{document}
\maketitle
\newpage
\footnotesize
\tableofcontents
\listoffigures
\newpage
\normalsize

\newpage
% Rolls {{{
\newpage
\section{Rouleaux difficiles d'accès}
\subsection{Description}

Toute la production de produits de vinyles repose sur le matériel principal,
le vinyle, qui arrive dans l'usine sous forme de rouleau. Le rangement de ces
rouleaux est pêle-mêle ce qui rends les rouleaux du bas et des côtés difficiles
d'accès sans causer de déboulements.

\subsection{Contexte}

Lors de la visite de la section des vinyles, ce coin du département est le
premier qui a été visité puisque cela consiste en la première étape pour la
production d'un produit finit.

\quickfig{Rangement des rouleaux à vinyle}{rolls}{.6}{figures/rolls-2.jpg}

\subsection{En quoi est-ce un détail?}

Ce problème pourrait sembler anodin, puisqu’il n’entrave pas directement le fonctionnement
des machines ou la qualité du produit final. Il concerne uniquement la logistique de
manutention, un aspect souvent négligé car en périphérie du processus de transformation
proprement dit.

\subsection{En quoi est-ce une observation méprisée?}

La disposition des rouleaux est perçue comme une contrainte tolérée par les employés,
suggérant qu’aucune solution n’a été priorisée. L’espace de stockage est considéré comme
acquis et peu flexible, donc l’enjeu est ignoré ou relégué à une simple question d’effort
physique à fournir.

\subsection{En quoi est-ce une observation singulière?}

Elle révèle un point d’ancrage du processus qui n’est pas automatisé ni mécanisé, et qui
dépend entièrement de l’organisation spatiale et de la manipulation humaine. Cela met en
évidence un goulet d’étranglement potentiel dès l’étape initiale, rarement abordé dans
l’analyse des flux de production.
% }}}
\newpage
% Templates {{{
\section{Gabarits de production}
\subsection{Description}

Dans l'usine Polytech, il y a plusieurs machines qui ont besoin de gabarits
afin de permettre la production de leurs différents produits en vinyle. Tous
les gabarits sont placés sur un rack avec plus d'une centaine de gabarits
différents qui semblent un peu pêle-mêle et difficiles à retrouver rapidement.
Les échaffaudages requièrent l'utilisation d'une échelle pour être capable
d'aller chercher certains gabarits, ce qui rend l'installation peu pratique
lors de changements à la production.

\subsection{Contexte}

Cette observation à été faite peu de temps après celle des rouleaux. Les deux
sections du département étant connexes et très rapprochées physiquement au fond
du département.

\quickfig{Rangement des gabarits}{templates}{.3}{figures/templates.jpg}

\subsection{En quoi est-ce un détail?}

La disposition désorganisée des gabarits semble secondaire, car elle n’interfère pas
immédiatement avec la production tant que les opérateurs savent où chercher. Le système
fonctionne tant bien que mal, ce qui peut masquer l’inefficacité sous-jacente.

\subsection{En quoi est-ce une observation méprisée?}

Le recours à une échelle et la perte de temps associée sont banalisés comme faisant
partie du quotidien. Le manque d’organisation est accepté comme norme, et aucune mesure
corrective apparente n’a été mise en place, signe que cette contrainte est tolérée
plutôt que remise en question.

\subsection{En quoi est-ce une observation singulière?}

Elle met en lumière la dépendance humaine dans un processus pourtant industrialisé.
Le contraste entre la précision nécessaire des gabarits et leur stockage rudimentaire
souligne une tension entre exigence technique et improvisation logistique rarement évoquée.

% }}}
\newpage
% Tape {{{
\section{Perte de travail à chaque mise en place}

\subsection{Description}

Chaque poste de production est situé autour d'une machine qui positionne tout
les travailleurs en cercle autour, chacun ayant son poste. Chaque station a des
papiers collants avec de l'écriture désignant comment aligner/couper/plier
certains morceaux de vinyles pour en faire un produit. À chaque fois qu'une
machine doit être réassignée à la production d'un produit différent, toutes ces
notes sont perdues. Elles doivent être arrachées afin de faire place à de
nouvelles notes.

\subsection{Contexte}

L'observation a été faite lors d'un moment où les travailleurs étaient en pause
me permettant d'observer les machines d'un peu plus près.

\quickfig{Autocollants devant être refaits à chaque mise en place}{tape}{.6}{figures/tape-close.jpg}

\subsection{En quoi est-ce un détail?}

Cette perte d’information peut sembler mineure, car elle n’interrompt pas directement la
production et les opérateurs semblent habitués à la réécriture manuelle des notes. Le
processus de remplacement des étiquettes est rapide, ce qui le rend presque invisible
dans le quotidien.

\subsection{En quoi est-ce une observation méprisée?}

La disparition systématique de ces connaissances tacites à chaque changement de production
indique un manque de valorisation du savoir opérateur. L'absence d’un support plus durable
ou standardisé démontre que cette mémoire de travail est vue comme jetable et sans importance.

\subsection{En quoi est-ce une observation singulière?}

Elle révèle une forme de documentation vivante mais éphémère, maintenue uniquement par les
gestes et la mémoire des travailleurs. Cette perte cyclique d’informations expose une faille
dans la capitalisation du savoir pratique, rarement considérée dans les analyses de performance.

% }}}
\newpage
% Communication {{{
\section{Communication difficile}
\subsection{Description}

Le département de vinyle est assez bruillant ce qui rends la communication
entre les travailleurs et avec les contremaîtres assez difficile en plus de
rendre l'expérience de travail inférieure à celle qu'elle pourrait être.

\subsection{Contexte}

La discussion qui suit s'est produite non pas durant un pause, mais pendant que
toutes les machines marchaient à plein régime dans un environnement à aire
ouverte où il est possible d'entendre les machines de tous les autres
département en plus de celles du département.

\boxfig{Discussion ayant mené à la réalisation du niveau de bruit}{noise}{
	Un fois arrivé à l'atelier de vinyle (les autre sections de l'usine nous
	ayant été montrées juste avant), la personne en charge de la visite nous
	laisses entre les mains de la personne en charge de la production de produits
	vinyles. Nous commençons à lui poser plusieurs questions face à son
	environnement. Malgré le fait qu'ils réponds à toutes nos questions, nous
	arrivons difficilement à comprendre tous ses propos vu la quantité de bruits
	présents dans l'usine.
}

\subsection{En quoi est-ce un détail?}

Le bruit est souvent perçu comme une composante normale du milieu industriel.
Il n'affecte pas directement la production matérielle, ce qui en fait un enjeu
facilement relégué au second plan malgré son impact humain et organisationnel.

\subsection{En quoi est-ce une observation méprisée?}

Les difficultés de communication ne sont pas activement adressées, malgré leur
influence sur la coordination et la sécurité. Le bruit est toléré comme une
contrainte inévitable, sans mesure apparente pour le réduire ou le contourner.

\subsection{En quoi est-ce une observation singulière?}

Elle met en lumière un facteur invisible mais omniprésent qui affecte
directement la qualité des interactions humaines. Ce bruit rend plus difficile
la transmission d’information, et donc la fluidité des opérations, ce qui
s’oppose au principe fondamental d’efficacité dans un environnement manufacturier.

% }}}
% Setup Time {{{
\newpage
\section{Temps de mise en place}
\subsection{Description}

Chaque machine agissant comme chaîne de production pour un produit de vinyle
nécessite une mise en place afin de permetre à un nouveau produit d'être
manufacturé (selon les demandes). Beaucoup de produits ne sont pas produits en
énormes quantité ce qui résulte en un temps de mise en place des machines qui
est plus long que le temps de production en soit.

\subsection{Contexte}

Voyant diverses machines semblant tous travailler sur des produits
différents, je remarques qu'une machine est en train d'être mise en place
pour la production d'un nouveau produit.

\boxfig{Discussion ayant dévoilé le temps de mise en place}{setup}{
	Je demande à celui en charge du département à quoi ressemble le
	processus de mise en place des machines et il me réponds que ce processus est
	assez long, souvent plus long que la production du produit qu'ils ont à
	faire ayant l'
}

\subsection{En quoi est-ce un détail?}

Ce phénomène peut sembler anecdotique lorsqu'on observe uniquement le produit final.
Le fait que le temps de mise en place dépasse celui de production n’est pas visible
extérieurement, surtout si le rendement global semble constant.

\subsection{En quoi est-ce une observation méprisée?}

Le temps de mise en place est traité comme un mal nécessaire, rarement optimisé ou
réduit de manière systémique. Il est intégré à la routine des opérateurs, mais sans
effort manifeste pour le minimiser, malgré son impact direct sur l’efficacité.

\subsection{En quoi est-ce une observation singulière?}

Elle révèle un paradoxe fondamental dans le modèle de production : une grande
flexibilité est exigée sans que l’infrastructure ne soit conçue pour des
changements fréquents. Cette tension entre variété et productivité expose une limite
structurelle du système.

% }}}

\end{document}
